% IEEE conference paper skeleton
% Advanced Telerobotics Research Lab, Department of Computer Science, Kent State University
%
% Build:  tectonic main.tex        (or: latexmk -pdf main.tex)
% Guide:  ../../../research/writing-papers.md
%
% Before you submit:
%   1. Check the venue's call for papers. Many IEEE Robotics and Automation
%      Society conferences (ICRA, IROS, Humanoids) supply their own template
%      (ieeeconf.cls); if so, move this body into that file.
%   2. Double-blind venue? Set \anonymoustrue below and remove self-identifying
%      text (acknowledgments, "our previous work [3]", repository links).
%   3. Replace every [bracketed placeholder]. Search the file for "[".

\documentclass[conference]{IEEEtran}
\IEEEoverridecommandlockouts % allows \thanks in the title block

% IEEE papers are set in Times. pdfLaTeX (Overleaf's default) gets Times from
% IEEEtran itself; XeLaTeX, LuaLaTeX and tectonic need a Times-compatible font.
\usepackage{iftex}
\ifPDFTeX\else
  \usepackage{fontspec}
  \setmainfont{texgyretermes}[Extension=.otf, UprightFont=*-regular,
    BoldFont=*-bold, ItalicFont=*-italic, BoldItalicFont=*-bolditalic]
\fi

\usepackage{cite}
\usepackage{amsmath,amssymb,amsfonts}
\usepackage{graphicx}
\usepackage{textcomp}
\usepackage{xcolor}
\usepackage{url}

% ---------------------------------------------------------------------------
% Review toggle
% ---------------------------------------------------------------------------
\newif\ifanonymous
\anonymousfalse % \anonymoustrue hides authors and acknowledgments for double-blind review

% ---------------------------------------------------------------------------
% Title and authors
% ---------------------------------------------------------------------------
\title{[Paper Title in the Venue's Capitalization Style]}

\ifanonymous
  \author{\IEEEauthorblockN{Anonymous Author(s)}
    \IEEEauthorblockA{Affiliation withheld for review}}
\else
  \author{%
    \IEEEauthorblockN{[First Author Name]}
    \IEEEauthorblockA{\textit{Advanced Telerobotics Research Lab}\\
      \textit{Department of Computer Science}\\
      \textit{Kent State University}\\
      Kent, OH, USA\\
      {[username]@kent.edu}}
    \and
    \IEEEauthorblockN{[Second Author Name]}
    \IEEEauthorblockA{\textit{Advanced Telerobotics Research Lab}\\
      \textit{Department of Computer Science}\\
      \textit{Kent State University}\\
      Kent, OH, USA\\
      {[username]@kent.edu}}
    \and
    \IEEEauthorblockN{[Faculty Advisor Name]}
    \IEEEauthorblockA{\textit{Advanced Telerobotics Research Lab}\\
      \textit{Department of Computer Science}\\
      \textit{Kent State University}\\
      Kent, OH, USA\\
      {[username]@kent.edu}}
  }
\fi

\begin{document}

\maketitle

\begin{abstract}
[One paragraph, usually 150 to 250 words; check the call for papers. State the
problem, what you built or studied, how you evaluated it and the main result
with its conditions (for example, ``in simulation'' or ``on a Unitree Go2 in the
lab'').]
\end{abstract}

\begin{IEEEkeywords}
[telerobotics], [human-robot interaction], [keyword three]
\end{IEEEkeywords}

% ---------------------------------------------------------------------------
\section{Introduction}
% ---------------------------------------------------------------------------
[Motivate the problem and state your contributions as a short list.]
Robot software in the lab builds on ROS~2~\cite{macenski2022ros2,ros2jazzydocs},
and many of our controllers are first tested in simulation~\cite{todorov2012mujoco}.

\section{Related Work}
[Group prior work by approach and say how yours differs.]

\section{Method}
[Describe the system. Define every symbol the first time it appears.]
\begin{equation}
  \mathbf{u}_t = K\,(\mathbf{x}^{*}_t - \mathbf{x}_t)
  \label{eq:control}
\end{equation}
Equation~\eqref{eq:control} is a placeholder.

\section{Experiments}
[Setup, conditions, participants (with IRB protocol number once approved),
metrics and results. Say whether each robot run was teleoperated, autonomous
or scripted.]

% Sample figure. Replace the \fbox placeholder with
%   \includegraphics[width=\linewidth]{figures/[file].pdf}
% Export plots as PDF (vector) and photos as PNG or JPEG at 300 dpi or more.
\begin{figure}[t]
  \centering
  \fbox{\parbox[c][1.4in][c]{0.92\linewidth}{\centering Figure placeholder}}
  \caption{[One-sentence caption that states what the figure shows and under
  which conditions, for example ``Teleoperated trials, $n=[N]$'']}
  \label{fig:overview}
\end{figure}

Fig.~\ref{fig:overview} shows [what the reader should notice].

\begin{table}[t]
  \caption{[Table Caption Above the Table, as IEEE Style Requires]}
  \label{tab:results}
  \centering
  \begin{tabular}{lcc}
    \hline
    Condition & [Metric 1] & [Metric 2] \\
    \hline
    {[Baseline]} & {[0.00]} & {[0.00]} \\
    {[Ours]}     & {[0.00]} & {[0.00]} \\
    \hline
  \end{tabular}
\end{table}

\section{Discussion and Limitations}
[What the results do and do not show.]

\section{Conclusion}
[Two or three sentences.]

\ifanonymous\else
\section*{Acknowledgment}
% Name funders exactly as the award record does, and only after the PI confirms
% which award funded this work.
This work was supported in part by [Sponsor] under [Grant No.].
% If you used generative AI to write or edit text, or to make figures or code,
% disclose it here as the IEEE policy on AI-generated content requires
% (see research/standards.md).
\fi

\bibliographystyle{IEEEtran}
\bibliography{references}

\end{document}
